\documentclass[12pt,a4paper]{article}
% ============================================================
%  Structural Persistence Series — Shared Preamble
% ============================================================
\usepackage{fontspec}
\setmainfont{Hiragino Mincho ProN}
\setsansfont{Hiragino Sans}
\setmonofont{Menlo}

\usepackage{iftex}
\ifXeTeX
  \XeTeXlinebreaklocale "ja"
  \XeTeXlinebreakskip=0pt plus 1pt minus 0.1pt
\fi

\usepackage[margin=22mm]{geometry}
\usepackage{amsmath,amssymb,amsthm}
\usepackage{xcolor}
\usepackage{graphicx}
\usepackage{hyperref}
\usepackage{booktabs}
\usepackage{fancyhdr}
% enumitem not available in this TeX installation

% --- Colours ------------------------------------------------
\definecolor{PaperAccent}{HTML}{1F4E79}
\definecolor{PaperGray}{HTML}{51606F}

\hypersetup{
  colorlinks=true,
  linkcolor=PaperAccent,
  urlcolor=PaperAccent,
  citecolor=PaperAccent,
  pdflang=ja
}
\urlstyle{same}

% --- Typography ---------------------------------------------
\setlength{\parindent}{1.5em}
\setlength{\parskip}{0pt}
\setlength{\emergencystretch}{3em}
\linespread{1.08}
\setlength{\abovedisplayskip}{8pt plus 2pt minus 4pt}
\setlength{\belowdisplayskip}{8pt plus 2pt minus 4pt}
\setlength{\abovedisplayshortskip}{6pt plus 2pt minus 2pt}
\setlength{\belowdisplayshortskip}{6pt plus 2pt minus 2pt}

% --- Section label names ------------------------------------
\renewcommand{\abstractname}{要旨}

% --- Theorem-like environments ------------------------------
\theoremstyle{plain}
\newtheorem{proposition}{命題}[section]
\newtheorem{lemma}[proposition]{補題}
\newtheorem{corollary}[proposition]{系}

\theoremstyle{definition}
\newtheorem{definition}{定義}[section]
\newtheorem{assumption}{条件}[section]

\theoremstyle{remark}
\newtheorem*{remark}{注}

% --- Running header -----------------------------------------
\pagestyle{fancy}
\fancyhf{}
\renewcommand{\headrulewidth}{0.4pt}
\fancyhead[L]{\small\textcolor{PaperGray}{\nouppercase{\leftmark}}}
\fancyhead[R]{\small\textcolor{PaperGray}{\thepage}}
\fancypagestyle{plain}{%
  \fancyhf{}
  \renewcommand{\headrulewidth}{0pt}
  \fancyfoot[C]{\small\textcolor{PaperGray}{\thepage}}
}

% --- Title block macro --------------------------------------
\newcommand{\PaperTitleBlock}[3]{%
  % #1 = title, #2 = subtitle, #3 = date
  \thispagestyle{plain}%
  \begin{center}
  {\LARGE\bfseries #1\par}
  \vspace{0.4em}
  {\normalsize #2\par}
  \vspace{1.0em}
  {\large Akihito Sunagawa\par}
  \vspace{0.3em}
  {\normalsize #3\par}
  \end{center}
  \vspace{1.6em}
}

% Legacy 2-arg version (backward compat)
\newcommand{\PaperTitleBlockSimple}[2]{%
  \PaperTitleBlock{#1}{}{#2}
}


\begin{document}

\PaperTitleBlock
  {持続知能の混成構成}
  {--- なぜ外部代謝が必要なのか ---}
  {2026年4月}

\begin{abstract}
持続的に運用される知能系は、新しい知識を学び、矛盾を処理し、過去の知識構造を壊しすぎずに維持しなければならない。本稿は、この課題に対して、パラメータ的適応だけでも、外部記憶だけでも十分ではないことを論じる。これまでの一連の結果は、構造持続の最小形式~\cite{paper1}、条件つきの指数導出~\cite{paper2}、論理一貫性維持の具体例~\cite{paper3}、そして前提認識型の継続学習実験~\cite{paper4}を通じて、知識持続の失敗が構造損失として理解できることを示してきた。本稿では、それらを統合し、持続知能には少なくとも三つの能力が必要であることを主張する。すなわち、(1)~新しい状況への迅速なパラメータ的適応、(2)~矛盾と更新を長期的に整理する外部代謝、(3)~取得・整理された知識を実際の応答に反映する応答生成である。本稿の主張は、知識持続の問題が単なる記憶容量の問題ではなく、構造損失、矛盾処理、応答生成の三層からなるという点にある。
\end{abstract}

% =============================================================
\section{はじめに}
% =============================================================

本稿では、新しい知識を受け入れ、既存知識との矛盾を整理し、なお一貫した振る舞いを長期的に維持する能力を、持続知能と呼ぶ。持続知能は、単に一度学習して終わるものではなく、長期的な運用のなかで知識構造の整合性を維持し続けることを要請する。

このとき問題になるのは、記憶容量そのものではない。重要なのは、更新が既存の知識構造をどのように壊し、またどのように維持しうるかである。すでに最小形式~\cite{paper1}で示したように、構造持続は、構造を維持できる状態の残存と、維持に利用できる資源の双方に依存する。さらに、条件つき導出~\cite{paper2}により、その指数型は仮定ではなく条件つきの数学的帰結として与えられることを示した。

しかし、知識持続を現実の知能系に適用するとき、問題は二重になる。第一に、パラメータ的更新は新しい知識を素早く取り込めるが、矛盾を伴う更新のもとではしばしば上書き的に働く。第二に、外部記憶は古い知識と新しい知識を併置しうるが、それだけではモデルがそれを正しく使うことを保証しない。ここで外部代謝とは、更新された知識と無効化された知識を時間的・条件的に整理し、矛盾対を保持したまま検索可能な形に再編成する機構を指す。

本稿の目的は、この二重性を統合的に記述することである。本稿は新しい理論や実験を提示する研究論文ではなく、これまでの結果を統合し、持続知能の設計課題を明示する統合覚書である。ここで主張するのは、持続知能には単一の解決策では足りず、パラメータ的適応と外部代謝を組み合わせた混成構成が必要であるということである。

% =============================================================
\section{構造持続から見た知識更新}
% =============================================================

構造持続の最小形式~\cite{paper1}では、制約の蓄積は構造を維持できる状態集合の縮小として書かれ、累積損失は指数型に現れる。この見方を知識更新に移すと、矛盾を含む更新は単なる新知識の追加ではなく、知識構造を保てる状態の領域を縮小させる作用として読める。

このとき、失敗には二つの側面がある。一つは資源側の不足であり、もう一つは知識構造そのものの縮小である。前者は処理能力、時間、探索余地などに対応し、後者は整合的に知識を維持できる可能性の減少に対応する。持続知能とは、この二つを同時に管理する能力にほかならない。

この見方は、すでに個別論文において具体化されている。論理一貫性維持~\cite{paper3}では、有効推論経路の集合の縮小として構造損失を記述した。前提更新依存的継続学習~\cite{paper4}では、T2 collapse は有効な知識状態の集合が急激に縮小する点として読まれ、依存認識型の再提示はその損失の蓄積を抑える介入として、多アダプタ分離は維持資源を別部分空間に確保しようとする試みとして読まれる。

% =============================================================
\section{パラメータ的適応の役割と限界}
% =============================================================

LoRA を用いた前提更新依存的継続学習の結果~\cite{paper4}は、パラメータ的更新の役割と限界を明瞭に示す。実際に、7モデル・4系統の全条件で、最初の前提更新以降に T2 collapse が例外なく観測された。単純な再提示は忘却を遅らせるが、矛盾を伴う更新では前提と派生知識の一貫性を十分に保てない。前提認識制御器は依存整合性を改善するが、依然として更新は上書き的である。多アダプタ分離は保持と更新の衝突を部分的に緩和するが、保持忠実度は低い。

このことは、パラメータ的適応が不要だということを意味しない。むしろ逆である。新しい状況に対して素早く局所的に適応する能力は、持続知能に不可欠である。しかし、その能力だけでは、高忠実度の長期保持と矛盾管理までは担えない。

% =============================================================
\section{外部代謝の役割}
% =============================================================

前稿~\cite{paper4}で示唆したように、パラメータ的な更新だけでは高忠実度の長期保持に限界がある。この点から自然に要請されるのが、外部的な知識構造の維持・再編成機構である。

外部代謝は、更新された知識と無効化された知識を時間的・条件的に整理し、矛盾の対を保持したまま、必要時に検索可能な形に再編成する。ここで重要なのは、外部記憶が単に保存場所を提供するだけでなく、知識を構造化し、矛盾対を固定し、後の検索可能性を維持する点である。

この観点から見ると、外部代謝はパラメータ的更新の補助ではなく、別種の機能である。パラメータ的更新が局所的な重み変化によって現在の行動能力を調整するのに対し、外部代謝は長期にわたる知識構造の再編と維持を担う。

% =============================================================
\section{応答生成の障壁}
% =============================================================

しかし、パラメータ的適応と外部代謝を組み合わせても、問題はなお残る。取得された知識が、モデルの応答生成に正しく反映されない場合があるからである。予備的な対話実験では、検索の成功率が 96\% であったにもかかわらず、事実再現率は 6.7\% にとどまる事例が観測された。モデルは正しい情報を受け取っていながら、対話履歴や応答習慣に引かれて検索結果を無視し、誤った出力を返していた。したがって、知識の保持と検索だけでは十分ではなく、取得された知識を応答へ忠実に反映する応答生成の層を、独立の障壁として扱う必要がある。

このことは、知識持続の失敗が三層に分かれることを示す。第一層は保存と更新である。第二層は構造化と検索である。第三層は、取得された知識を実際の応答へと変換する応答生成である。

もし検索が成功しても、モデルがその内容を適切に用いないなら、知識持続は表面的には失敗として現れる。したがって、持続知能の問題は保存や矛盾管理だけでは完結しない。応答生成そのものが第三の障壁となる。

% =============================================================
\section{なぜ混成構成が必要か}
% =============================================================

以上をまとめると、持続知能には少なくとも三つの能力が必要である。

\begin{enumerate}
  \item \textbf{パラメータ的適応.} 新しい課題や局所的更新に素早く反応するために必要である。
  \item \textbf{外部代謝.} 矛盾する知識を整理し、長期的な知識構造を維持するために必要である。
  \item \textbf{応答生成の改善.} 検索された知識や整理済みの知識構造を、実際の出力へ忠実に反映するために必要である。
\end{enumerate}

この三者は相互に代替可能ではない。パラメータ的適応だけでは高忠実度の長期保持が難しい。外部記憶だけでは局所的適応と出力反映が不足する。応答生成の改善だけでは、そもそも保持されていない知識を救えない。

したがって、持続知能は単一の工夫によって達成されるものではなく、三層を統合した混成構成を必要とする。

% =============================================================
\section{位置づけ}
% =============================================================

本稿は、最小形式の一般論を拡張するものではない。むしろ、それまでの結果を統合し、持続知能という設計課題の上に置き直すことを目的とする。

最小形式~\cite{paper1}は、構造損失を表す骨格を与える。条件つき導出~\cite{paper2}は、その骨格の数学的基盤を与える。論理一貫性維持~\cite{paper3}は、観測量の悪化と構造損失の接続を最初の具体例として示す。前提更新依存的継続学習~\cite{paper4}は、パラメータ的適応の上書き的限界と、その部分的緩和を示す。本稿は、それらをまとめ、持続知能のための設計原理として読み直す位置にある。この方向の一つの具体化として、外部記憶と代謝パイプラインを組み合わせた対話系の試作が行われているが、その評価と詳細は別稿に譲る。

% =============================================================
\section{限界}
% =============================================================

本稿は、完全な混成知能アーキテクチャを実装したことを主張するものではない。また、三つの能力の最適な統合方法をすでに与えたわけでもない。

本稿が与えるのは、より限定された主張である。すなわち、これまでの理論と実験を総合すると、持続知能は単なるパラメータ的学習の延長としては捉えきれず、外部代謝と応答生成の問題を同時に含む、ということである。なお、三層間の競合——たとえばパラメータ的適応が更新した方向と外部代謝が保持すべきと判断した方向が衝突する場合——の解決方針は、本稿では扱わない。

% =============================================================
\section{結論}
% =============================================================

持続知能の課題は、知識を保存することそのものではない。更新、矛盾、再編成、検索、そして応答への反映を、長期的に整合したかたちで維持することである。

本稿では、これまでの結果を統合し、持続知能には少なくとも三つの能力が必要であることを論じた。すなわち、パラメータ的適応、外部代謝、応答生成の改善である。結論として、持続知能は単なる重み更新でも、単なる外部記憶でも足りない。必要なのは、それらを構造損失の観点から統合した混成構成である。

本稿の役割は、その必要性を、構造持続という共通の座標のもとで明示することにある。

% =============================================================
\begin{thebibliography}{9}
\bibitem{paper1} Sunagawa, A. 構造持続の最小形式. 2026.
\bibitem{paper2} Sunagawa, A. 構造持続の条件つき導出. 2026.
\bibitem{paper3} Sunagawa, A. 論理一貫性維持の構造持続的記述. 2026.
\bibitem{paper4} Sunagawa, A. 前提更新を伴う継続学習における構造的忘却. 2026.
\end{thebibliography}

\end{document}
